\documentclass[UTF8]{ctexart}
\usepackage{xeCJK}
\usepackage{geometry}
\geometry{left=2cm,right=2cm,top=2.5cm,bottom=2.5cm}
\setlength{\headheight}{12.7pt}

\usepackage{titlesec}
\titleformat{\section}
  {\normalfont\large\bfseries}
  {\thesection}
  {1em}
  {}
\titlespacing*{\section}
  {0pt}
  {3.5ex plus 1ex minus .2ex}
  {2.3ex plus .2ex}

\usepackage{amsmath}
\usepackage{amssymb}
\usepackage{amsthm}
\usepackage{enumitem}
\setlist[enumerate]{itemsep=0pt,topsep=0pt,parsep=0pt,leftmargin=0pt,itemindent=2.5em}
\setlist[itemize]{itemsep=0pt,topsep=0pt,parsep=0pt,leftmargin=0pt,itemindent=2.5em}
\usepackage{fancyhdr}
\pagestyle{fancy}
\fancyhf{}
\usepackage{graphicx}
\usepackage{hyperref}
\usepackage{indentfirst}
\usepackage{lmodern}


\begin{document}
\setlength{\abovedisplayskip}{1pt}
\setlength{\belowdisplayskip}{1pt}

\section{全序关系}

\hypertarget{sec:定义1.1}{}
\noindent\textbf{定义1.1 (全序关系)}

给定集合$a$和$a$上的二元关系$\leqslant$. 如果以下四条成立:
\begin{enumerate}
    \item 自反性: $\forall x\in a:\,x\leqslant x$,
    \item 传递性: $\forall x,y,z\in a:\,
    x\leqslant y\land y\leqslant z\to x\leqslant z$,
    \item 反对称性: $\forall x,y\in a:\,x\leqslant y\land y\leqslant x\to x=y$,
    \item 可比性: $\forall x,y\in a:\,x\leqslant y\lor y\leqslant x$.
\end{enumerate}
那么称$\leqslant$是$a$上的一个\textbf{全序关系}.
也称$(a,\leqslant)$是一个\textbf{全序结构}.

\vspace{10pt}

给定全序结构$(a,\leqslant)$, 它衍生的$<$满足以下三条, 由此称之为\textbf{严格全序}:
\begin{enumerate}
    \item 反自反性: $\forall x\in a:\,x\nless x$,
    \item 传递性: $\forall x,y,z\in a:\,x<y\land y<z\to x<z$,
    \item 三歧性: $\forall x,y\in a:\,x<y\lor x=y\lor x>y$,
\end{enumerate}
另外, 此时$\geqslant$也是全序, $>$也是严格全序.

\vspace{10pt}

在全序中, 证明一个映射是序同构只需证明它单向的保序性质.

\vspace{10pt}

\hypertarget{sec:定理1.1}{}
\noindent\textbf{定理1.1}

给定全序结构$(a,\leqslant_a)$和$(b,\leqslant_b)$, 如果有满射$f:a\to b$使得
\[
    \forall x,y\in a:\,\,x<_ay\to f(x)<_bf(y),
\]
那么就有$f:(a,\leqslant_a)\cong(b,\leqslant_b)$.

\noindent{\kaishu 证明.}

对任意的$x,y\in a$, 首先显然有$x\leqslant_ay\to f(x)\leqslant_bf(y)$,
其次如果$f(x)\leqslant_bf(y)$但是$y<_ax$, 那么
\[
    y<_ax\implies f(y)<_bf(x)\implies\bot,
\]
从而$x\leqslant_ay$.

最后, 如果$f(x)=f(y)$, 那么同时有$x\leqslant_ay$和$y\leqslant_ax$,
即$x=y$. 所以$f$是单射.

综上即得$f:(a,\leqslant_a)\cong(b,\leqslant_b)$. {\hfill $\qedsymbol$}

\vspace{10pt}

\hypertarget{sec:定理1.2}{}
\noindent\textbf{定理1.2}

非空有限集上的全序一定有最小(大)元.

\noindent{\kaishu 证明.}

非空有限集上的全序有极小元, 极小元就是最小元. 最大元同理. \hfill $\qedsymbol$





\newpage

\hypertarget{sec:定义1.2}{}
\noindent\textbf{定义1.2 (前段)}

给定全序结构$(a,\leqslant)$和$s\subset a$.
如果$\forall x\in s\,\forall y\in a:\,y<x\to y\in s$,
就称$s$为$a$的一个\textbf{前段}.

如果$s$是$a$的前段且$s\neq a$, 那么称$s$为\textbf{真前段}.

\vspace{10pt}

\hypertarget{sec:例1.1}{}
\noindent\textbf{例1.1}

给定全序结构$(a,\leqslant)$和$t\in a$, 定义
\[
    a_t\overset{\mathrm{def}}{=}\{x\in a\mid x<t\},
\]
则$a_t$是$a$的真前段.

\noindent{\kaishu 证明.}

对任意的$x\in a_t$和$y<x$, 有
\[
    y<x\land x<t\implies y<t\implies y\in a_t,
\]
所以$a_t$是前段. 并且显然$t\notin a_t$, 所以$a_t$是真前段. \hfill $\qedsymbol$

\vspace{10pt}

易证对任意的$t,t'\in a$, $t\leqslant t'$当且仅当$a_t\subset a_{t'}$.

\vspace{10pt}

\hypertarget{sec:定理1.3}{}
\noindent\textbf{定理1.3}

任给全序的序同构$f:(a,\leqslant_a)\cong(b,\leqslant_b)$
和$a$的前段$s$, $f[s]$也是$b$的前段.

进一步, 如果存在$t\in a$使得$s=a_t$, 那么$f[s]=b_{f(t)}$.

\noindent{\kaishu 证明.}

对任意的$x\in f[s]$和$y<_bx$, 有
\[
    f^{-1}(x)\in s\land f^{-1}(y)<_{a} f^{-1}(x)\implies
    f^{-1}(y)\in s\implies y\in f[s],
\]
所以$f[s]$也是前段.

进一步, 如果存在$t\in a$使得$s=a_t$, 那么对任意的$y$有
\begin{equation*}
    y\in f[s]\iff f^{-1}(y)\in s\iff f^{-1}(y)<_at
    \iff y<_bf(t)\iff y\in b_{f(t)}.
\tag*{\qedsymbol}\end{equation*}

% 今天写的同样很多, 虽然大部分是复制过来的...
% Can you feel the music ?
% 2026.04.10

\vspace{10pt}

\hypertarget{sec:定理1.4}{}
\noindent\textbf{定理1.4}

任给全序结构$(a,\leqslant)$, 记$S$为其全体前段的集合, 则$(S,\subset)$也是全序结构.

\noindent{\kaishu 证明.}

对任意两个前段$s,s'$, 假设$s\not\subset s'$且$s'\not\subset s$,
可以取$x\in s\setminus s'$, $y\in s'\setminus s$.

因为$a$是全序, 所以:
\begin{enumerate}
    \item 如果$x<y$, 那么由$s'$是前段得$x\in s'$, 矛盾;
    \item 如果$x=y$, 显然矛盾;
    \item 如果$x>y$, 那么由$s$是前段得$y\in s$, 矛盾. \hfill $\qedsymbol$
\end{enumerate}

\end{document}